% Work-MAT | 2. rocnik | hodina 15
\documentclass[11pt,a4paper]{article}
\input{../../sablony/pripravy/priprava.sty}
\nastavPripravu{2. ročník}{Exponenciální funkce}{15}{Exponenciální závislosti a modely}
\begin{document}
\poznamka{Druhá část „exponenciální funkce se vším všudy“: matematický model si žáci sestavují sami ze slovní situace. Nevyužívej logaritmy; hledáme hodnoty v dopředu zadaných časových okamžicích, nikoli obecný čas dosažení prahu.}
\blok{Opakované přičítání, nebo opakované násobení?}
Ne každá závislost v čase je exponenciální. Stálý absolutní přírůstek vede k lineárnímu modelu, zatímco stálý násobný koeficient k exponenciálnímu. U procentního růstu se nové procento vždy počítá \emph{z nové, již změněné hodnoty}.
\begin{prikladReseniBox}[title={Příklad 1}]
Objem kapaliny je na začátku 200 litrů a každou hodinu vzroste o $10\,\%$ aktuálního objemu. Sestav model a spočti objem po třech hodinách.
\tcblower
Koeficient za den je $a=1+0{,}10=1{,}10$. Proto
\[
N(t)=200\cdot(1{,}1)^t,\qquad N(3)=200\cdot1{,}331=266{,}2.
\]
Matematický model dává po třech hodinách objem 266{,}2 litru.
\end{prikladReseniBox}
\begin{prikladReseniBox}[title={Příklad 2}]
Hodnota stroje je 80\,000 Kč a každoročně klesne o $15\,\%$ své aktuální hodnoty. Zapiš model a vypočti hodnotu po dvou letech.
\tcblower
Každý rok zůstane $85\,\%$ předchozí hodnoty, tedy $a=0{,}85$.
\[
V(t)=80\,000\cdot(0{,}85)^t,\qquad V(2)=80\,000\cdot0{,}7225=57\,800\text{ Kč}.
\]
\end{prikladReseniBox}
\blok{Sestavení modelu z tabulky}
U dvou po sobě jdoucích hodnot při stejné délce časového kroku určíme koeficient poměrem $a=\frac{N(t+1)}{N(t)}$, pokud je jmenovatel nenulový. U skutečné exponenciální závislosti tento poměr zůstává stejný.
\begin{prikladReseniBox}[title={Příklad 3}]
V časech $t=0,1,2,3$ byly naměřeny hodnoty $160,120,90,67{,}5$. Rozhodni, zda jde o exponenciální model, a sestav jeho předpis.
\tcblower
\[
\frac{120}{160}=\frac{90}{120}=\frac{67{,}5}{90}=\frac{3}{4}.
\]
Poměr je konstantní. Počáteční hodnota je 160, proto $N(t)=160\cdot\left(\frac{3}{4}\right)^t$.
\end{prikladReseniBox}
\begin{prikladReseniBox}[title={Příklad 4}]
Základní částka 5\,000 Kč se úročí jednou ročně sazbou $4\,\%$. Sestav model konečné částky a vypočti stav po třech letech bez dalších vkladů.
\tcblower
\[
S(n)=5\,000\cdot(1{,}04)^n,\qquad S(3)=5\,624{,}32\text{ Kč}.
\]
Jde o jednoduchý model ročního složeného úročení bez daní a poplatků.
\end{prikladReseniBox}
\poznamka{Ve školní diskusi se osvědčuje otázka: „Přidáme každý rok 100 korun, nebo každý rok 5 procent?“ Odpověď má žák odvodit ze slovního zadání, ne jen odhadnout podle tvaru grafu.}
\begin{zadaniReseniBox}[title={Samostatná práce 1}]
\begin{enumerate}[label=\alph*)]
\item Populace 4\,000 bakterií se každou hodinu zdvojnásobí. Zapiš model a stav po třech hodinách.
\item Zařízení si každý rok zachová $90\,\%$ hodnoty. Počáteční hodnota je 12\,000 Kč. Zapiš hodnotu po dvou letech.
\item Pro posloupnost $50,60,72,86{,}4$ urči násobný koeficient a exponenciální předpis.
\end{enumerate}
\tcblower
\begin{enumerate}[label=\alph*)]
\item $N(t)=4\,000\cdot2^t$, $N(3)=32\,000$.
\item $V(t)=12\,000\cdot(0{,}9)^t$, $V(2)=9\,720$ Kč.
\item $a=\frac{6}{5}=1{,}2$, $f(n)=50\cdot(1{,}2)^n$.
\end{enumerate}
\end{zadaniReseniBox}
\begin{zadaniReseniBox}[title={Rozšíření 1}]
Laboratorní vzorek má na začátku hmotnost 64 g. Každé čtyři hodiny zbývá polovina aktuální hmotnosti. Sestav model pro čas $t$ v hodinách a vypočti hmotnost po dvanácti hodinách.
\tcblower
Během $t$ hodin proběhne $\frac{t}{4}$ čtyřhodinových intervalů. Proto $m(t)=64\cdot\left(\frac{1}{2}\right)^{\frac{t}{4}}$. Pro $t=12$ je $m(12)=64\cdot\left(\frac{1}{2}\right)^3=8$ g.
\end{zadaniReseniBox}
\begin{zadaniReseniBox}[title={Domácí úkol}]
\begin{enumerate}[label=\alph*)]
\item Výrobek s cenou 2\,500 Kč zdraží za každý rok o $6\,\%$. Sestav model a urč cenu za dva roky.
\item Ze 400 g látky zbývá na konci každé hodiny $80\,\%$. Urči hodnotu za tři hodiny.
\item Vysvětli, proč řada $100,110,120,130$ nevyjadřuje stálý desetiprocentní růst.
\end{enumerate}
\tcblower
\begin{enumerate}[label=\alph*)]
\item $C(t)=2\,500\cdot(1{,}06)^t$, $C(2)=2\,809$ Kč.
\item $m(t)=400\cdot(0{,}8)^t$, $m(3)=204{,}8$ g.
\item Každý krok přidává 10 kusů, zatímco poměr dvou sousedních hodnot se mění.
\end{enumerate}
\end{zadaniReseniBox}
\end{document}
